\documentclass[a4paper,10pt]{article}
\usepackage[czech]{babel}
\usepackage[T1]{fontenc}
\usepackage[utf8]{inputenc}
\usepackage{amsmath,amssymb}
\usepackage[margin=18mm]{geometry}
\setlength{\parindent}{0pt}
\begin{document}
\section*{Exponenciální rovnice A/B -- řešení a bodování}
Podklad pro učitele k testu:\\
\texttt{01\_test\_ExponencialniRovnice\_AB.tex}.

Zařazení: 2. ročník SPŠ, plán \texttt{2\_rocnik/tematicke\_plany/SPS.pdf},
s.~1, listopad, bod 1.7. Společná listopadová dotace pro exponenciální
rovnice a nerovnice je 12 hodin. Test obsahuje pouze rovnice bez logaritmů.

\subsection*{Varianta A}
\begin{enumerate}
\item $D=\mathbb{R}$. $4^{x-1}=8 \iff 2^{2x-2}=2^3 \iff 2x-2=3$.
Tedy $K=\{\frac52\}$. Zkouška: $4^{3/2}=8$.
\item $D=\mathbb{R}$. $3^{x+1}+3^x=36 \iff 3^x(3+1)=36 \iff 3^x=9$.
Tedy $K=\{2\}$. Zkouška: $27+9=36$.
\item $D=\mathbb{R}$. Substituce $t=2^x>0$ dává
$t^2-5t+4=0 \iff (t-1)(t-4)=0$. Oba kořeny $t=1,4$ jsou kladné.
Z $2^x=1$ nebo $2^x=4$ plyne $K=\{0,2\}$.
Zkoušky: $1-5+4=0$, $16-20+4=0$.
\item $2^x-2\ne0 \iff x\ne1$, tedy $D=\mathbb{R}\setminus\{1\}$.
Pro $t=2^x>0$, $t\ne2$ platí
\[
\frac{t^2-4}{t-2}=6
\iff \frac{(t-2)(t+2)}{t-2}=6
\iff t+2=6 \iff t=4.
\]
Proto $x=2$ vyhovuje podmínce a $K=\{2\}$.
Zkouška: $(16-4)/(4-2)=6$. Hodnota $x=1$ je nepřípustná.
\end{enumerate}

\subsection*{Varianta B}
\begin{enumerate}
\item $D=\mathbb{R}$. $9^{x+1}=27 \iff 3^{2x+2}=3^3 \iff 2x+2=3$.
Tedy $K=\{\frac12\}$. Zkouška: $9^{3/2}=27$.
\item $D=\mathbb{R}$. $2^{x+2}+2^x=40 \iff 2^x(4+1)=40 \iff 2^x=8$.
Tedy $K=\{3\}$. Zkouška: $32+8=40$.
\item $D=\mathbb{R}$. Substituce $t=3^x>0$ dává
$t^2-10t+9=0 \iff (t-1)(t-9)=0$. Oba kořeny $t=1,9$ jsou kladné.
Z $3^x=1$ nebo $3^x=9$ plyne $K=\{0,2\}$.
Zkoušky: $1-10+9=0$, $81-90+9=0$.
\item $3^x-3\ne0 \iff x\ne1$, tedy $D=\mathbb{R}\setminus\{1\}$.
Pro $t=3^x>0$, $t\ne3$ platí
\[
\frac{t^2-9}{t-3}=12
\iff \frac{(t-3)(t+3)}{t-3}=12
\iff t+3=12 \iff t=9.
\]
Proto $x=2$ vyhovuje podmínce a $K=\{2\}$.
Zkouška: $(81-9)/(9-3)=12$. Hodnota $x=1$ je nepřípustná.
\end{enumerate}

\subsection*{Společné bodování a kontrola}
\begin{enumerate}
\item 2 body: převod na společný základ (1), vyřešení rovnice exponentů a výsledek (1).
\item 3 body: vytknutí (1), osamostatnění mocniny (1), výsledek (1).
\item 4 body: substituce s podmínkou $t>0$ (1), řešení kvadratické rovnice (1),
oba zpětné převody na $x$ (po 1).
\item 5 bodů: definiční podmínka (1), rozklad čitatele (1), korektní krácení (1),
vyřešení zbývající rovnice (1), přípustný výsledek (1).
\end{enumerate}
Celkem 14 bodů v každé variantě. Uznávejte i jiné správné postupy;
tutéž početní chybu netrestejte opakovaně, je-li další postup správný.
Orientační časy úloh: 2--3, 3--4, 4--5 a 5--6 minut,
celkem 14--18 minut a krátká kontrola. Odhad nebyl ověřen ve třídě.

Úplnost řešení plyne z ekvivalentních úprav, rozkladů kvadratických polynomů
a prostoty exponenciálních funkcí. Krácení v úlohách 4 platí na uvedeném
 definičním oboru. Varianty mají shodné bodování a srovnatelný počet kroků.
Všechny kořeny byly navíc nezávisle ověřeny výpočtem po dosazení do původních
rovnic. U úloh 4 byl ověřen i vyloučený kořen $x=1$ rovnice vzniklé
vynásobením jmenovatelem: odpovídající polynomy jsou $(t-2)(t-4)$
a $(t-3)(t-9)$.
\end{document}
